% ==========================================
% BAB I PENDAHULUAN
% ==========================================
\chapter{PENDAHULUAN}
\label{chap:pendahuluan}
% --- Latar Belakang ---
\section{Latar Belakang}
\label{sec:latar_belakang}

Pergeseran paradigma pendidikan tinggi global saat ini bergerak signifikan dari
pendekatan pedagogis yang berpusat pada pengajar menuju pendekatan yang berpusat
pada siswa \autocite{Ryoo2021}. Transformasi ini didorong oleh perkembangan
\textit{Personalized Learning} (PL), sebuah metode pendidikan yang menyesuaikan
materi, urutan, dan kecepatan belajar berdasarkan keterampilan serta preferensi
unik setiap individu. Kunci utama dalam merealisasikan personalisasi ini adalah
pemanfaatan kecerdasan artifisial (AI) melalui sistem \textit{Adaptive
Assessment} (AA), yang mampu secara dinamis menyesuaikan tingkat kesulitan konten
berdasarkan kinerja dan pengetahuan yang ditunjukkan oleh peserta didik secara
\textit{real-time} \autocite{Halkiopoulos2024}.

Menyadari potensi transformatif ini, pemerintah Indonesia telah menetapkan
pendidikan dan riset sebagai salah satu dari lima bidang prioritas dalam Strategi
Nasional Kecerdasan Artifisial (Stranas KA) 2020--2045 \autocite{BPPT2020}.
Dokumen strategis ini secara eksplisit menyerukan ditinggalkannya pendekatan
kaku ``\textit{one size fits all}'' dan mendorong eksplorasi pengembangan sistem
penilaian cerdas yang adaptif. Secara teknis, berbagai algoritma telah
dikembangkan untuk mendukung visi ini, mulai dari \textit{Bayesian Knowledge
Tracing} (BKT) yang memodelkan transisi penguasaan keterampilan, hingga
\textit{Deep Knowledge Tracing} (DKT) yang memanfaatkan \textit{Recurrent Neural
Network} (RNN) untuk pelacakan pengetahuan yang lebih presisi \autocite{Minn2022}.
Metode yang lebih praktis dan dapat diinterpretasikan seperti sistem
\textit{Elo rating} juga telah terbukti efisien secara komputasi untuk
mengestimasi kemampuan siswa dan kesulitan soal secara dinamis
\autocite{Vesin2022}.

Meskipun teknik-teknik tersebut unggul dalam mengoptimalkan jalur
belajar individu, penerapan yang naif terhadap algoritma yang hanya berfokus
pada preferensi historis individu membawa risiko laten yang serius. Dalam
literatur sistem rekomendasi modern, fenomena ini dikenal sebagai masalah
\textit{filter bubbles}, dalam fenomena ini algoritma secara tidak sengaja
mengisolasi
pengguna dari keragaman informasi dan perspektif dengan hanya menyajikan konten
yang selaras dengan riwayat mereka \autocite{Wang2022}. Dalam konteks pendidikan,
risiko ini bermanifestasi sebagai \textit{overpersonalization} (personalisasi
berlebih), yang menyebabkan sistem terlalu agresif dalam menyesuaikan konten dengan
kenyamanan individu sehingga justru mengasingkan siswa dari manfaat pembelajaran komunitas
\autocite{Halkiopoulos2024}. Algoritma yang bekerja semata-mata berdasarkan
prinsip homofili (kemiripan) cenderung menciptakan \textit{echo chambers}
akademik \autocite{Tommasel2021}, menghalangi interaksi sosial yang krusial bagi
pengembangan keterampilan kolaboratif.

Tinjauan literatur terkini mengonfirmasi adanya kesenjangan teknologi dalam
mengatasi dilema ini. Mayoritas sistem asesmen adaptif yang ada saat ini masih
berfokus eksklusif pada pelacakan pengetahuan individu tanpa mengintegrasikan
mekanisme penilaian sosial \autocite{Ihichr2024}. Model-model yang berfokus
tunggal seperti DKT atau sistem Elo standar secara inheren sulit menyeimbangkan
antara memberikan materi yang paling sesuai kemampuan individu dengan materi
yang memicu pertumbuhan melalui interaksi kelompok \autocite{Dokoupil2022}.

Kondisi ini menghadirkan tantangan sekaligus peluang strategis bagi implementasi
Stranas KA di Indonesia. Jika pengembangan sistem pendidikan nasional hanya
mengadopsi teknologi adaptif konvensional, Indonesia berisiko mewarisi kelemahan
fundamental berupa isolasi pembelajaran tersebut. Oleh karena itu, pengembangan
sistem asesmen di Indonesia tidak seharusnya sekadar meniru, melainkan perlu
melakukan lompatan teknologi dengan merancang sistem yang sejak awal memitigasi
risiko \textit{overpersonalization}. Dibutuhkan sebuah pendekatan baru yang
mampu mengelola keseimbangan antara adaptivitas personal dan dinamika kolaboratif,
memastikan bahwa intervensi teknologi memperkuat, bukan melemahkan, ekosistem
pembelajaran sosial.

% --- Rumusan Masalah ---
\section{Rumusan Masalah}
\label{sec:rumusan_masalah}
Berdasarkan latar belakang tersebut, teridentifikasi adanya kebutuhan mendesak
untuk mengembangkan sistem asesmen yang selaras dengan Stranas KA, namun tahan
terhadap risiko isolasi pembelajaran yang kerap muncul pada sistem adaptif
konvensional. Ketiadaan sistem yang mampu menyeimbangkan adaptivitas personal
dengan interaksi sosial menjadi kesenjangan teknologi utama.

Maka dari itu, permasalahan utama yang ingin diselesaikan adalah:
\begin{enumerate}
  \item Implementasi purwarupa sistem
  asesmen adaptif seperti apa yang mampu memitigasi efek
  \textit{overpersonalization} dengan mengintegrasikan metrik penilaian
  individual ke dalam mekanisme pembelajaran kolaboratif?
  \item Bagaimana efektivitas purwarupa sistem dalam mempertahankan fungsionalitas adaptif (personalisasi) sekaligus memfasilitasi interaksi
  sosial yang relevan saat diujicobakan pada pengguna?
\end{enumerate}


% --- Tujuan ---
\section{Tujuan}
\label{sec:tujuan}
Tujuan utama dari tugas akhir ini adalah \textbf{membangun purwarupa sistem
asesmen adaptif kolaboratif yang dirancang untuk meminimalkan risiko
\textit{over-
personalization} dalam lingkungan pembelajaran daring}.

Secara spesifik, tujuan penelitian ini adalah:
\begin{enumerate}
  \item Implementasi sistem terintegrasi antara algoritma asesmen adaptif
  dengan fitur kolaboratif sehingga
  dapat merekomendasikan interaksi sosial yang relevan untuk
  mengatasi \textit{overpersonalization} dalam pembelajaran.
  \item Mengevaluasi purwarupa sistem melalui uji coba pengguna untuk mengukur
  keberhasilan sistem dalam memberikan pengalaman belajar yang adaptif tanpa
  mengorbankan aspek kolaborasi.
\end{enumerate}

% --- Batasan Masalah ---
\section{Batasan Masalah}
\label{sec:batasan_masalah}
Untuk memastikan penelitian ini tetap fokus dan dapat diselesaikan dalam jangka waktu yang ditentukan, maka ditetapkan batasan-batasan sebagai berikut:
\begin{enumerate}
  \item Permasalahan asesmen adaptif yang
  dikaji akan difokuskan pada konteks yang selaras dengan visi Strategi Nasional
  Kecerdasan Artifisial (Stranas KA) Indonesia. Ini berarti penelitian berfokus
  pada pengembangan sistem yang mendukung ekosistem kolaboratif, yang disesuaikan
  dengan kebutuhan dan potensi implementasi di lingkungan pendidikan tinggi di
  Indonesia.
  \item Fokus masalah utama penelitian ini adalah pada kesenjangan integrasi
  algoritmik antara asesmen adaptif individual dan asesmen kolaboratif. Penelitian
  tidak akan menginvestigasi masalah lain dalam asesmen adaptif (seperti deteksi
  gaya belajar atau deteksi kecurangan dalam asesmen), melainkan terbatas pada
  perancangan model yang menggabungkan metrik kemahiran personal dengan metrik
  kinerja sosial.
  \item Permasalahan yang dikaji terbatas pada penilaian luaran kognitif
  (\textit{cognitive outcomes}), yaitu pelacakan "keterampilan" dan "pengetahuan"
  siswa. Penelitian ini tidak akan mencakup perancangan model untuk menilai luaran
  afektif (seperti motivasi atau kepuasan) atau luaran perilaku (seperti durasi
  belajar) sebagai ukuran utama kemahiran.
  \item Implementasi purwarupa sistem akan diujicobakan pada populasi sampel terbatas.
  Merujuk pada konteks pendidikan tinggi di Indonesia dan ketersediaan sumber daya,
  uji coba akan difokuskan pada sekelompok mahasiswa dari angkatan tertentu
  di institusi tertentu.
\end{enumerate}
% Tuliskan batasan-batasan yang diambil dalam pelaksanaan tugas akhir. Batasan ini dapat dihindari (bersifat opsional, tidak perlu ada) jika topik atau judul tugas akhir dibuat cukup spesifik.

% --- Metodologi Pengerjaan TA ---
\section{Metodologi}

Pelaksanaan tugas akhir ini akan mengadopsi dan mengadaptasi metodologi yang digunakan oleh \textcite{Vesin2022} dalam penelitian mereka mengenai asesmen adaptif dan rekomendasi konten. Metodologi ini terdiri atas beberapa tahapan utama yang sistematis, dimulai dari studi pendahuluan hingga analisis data sebagai .

\begin{enumerate}
    \item \textbf{Studi Literatur:} Melakukan tinjauan sistematis terhadap penelitian terkait asesmen adaptif, pembelajaran kolaboratif, dan implementasi sistem serupa.
    
    \item \textbf{Perancangan dan Pengembangan Sistem:} Merancang arsitektur perangkat lunak serta mengimplementasikan modul-modul sistem, termasuk integrasi algoritma asesmen adaptif dan antarmuka pengguna kolaboratif.
    
    \item \textbf{Sesi Pengguna (Pengujian dan Evaluasi):} Melaksanakan studi dengan partisipan sebagai sampel populasi pengguna untuk mengumpulkan data kuantitatif dan kualitatif mengenai penggunaan dan pengalaman mereka terhadap purwarupa sistem. Data yang dikumpulkan dapat mencakup data interaksi (\textit{click-streams}, \textit{log files}), hasil asesmen (peringkat yang dihasilkan), serta umpan balik pengguna (misalnya melalui survei atau wawancara).
    
    \item \textbf{Analisis Data:} Menganalisis data yang terkumpul untuk mengevaluasi fungsionalitas purwarupa dan efektivitas sistem dalam konteks asesmen adaptif dan kolaboratif, guna menjawab rumusan masalah serta tujuan penelitian.
    
    \item \textbf{Penarikan Kesimpulan dan Pelaporan:} Menyusun laporan tugas akhir berdasarkan hasil analisis dan temuan penelitian.
\end{enumerate}

% \begin{enumerate}
% \item	Tahapan investigasi pengumpulan fakta di latar belakang untuk merumuskan masalah.
% \item	Langkah-langkah pencarian, pengelompokan, dan penapisan literatur atau sumber informasi untuk mengumpulkan informasi yang relevan tentang topik yang diangkat, termasuk teori (konsep atau teori apa saja yang perlu dicari), hal-hal yang telah dicapai oleh orang lain (cara mencari dan kata kuncinya), dan berbagai informasi pendukung, untuk mencari solusi terhadap masalah yang dibahas. Gunakan metodologi yang tepat dalam menggali informasi dan dokumentasikan prosesnya (termasuk rekaman wawancara atau survei) di dalam Lampiran, termasuk tautan ke video atau foto. Hasil penggalian informasi ini akan dijelaskan secara sistematis di Bab II Studi Literatur.
% \end{enumerate}